\section{问题分析}

\subsection{总体分析}

本题的核心困难并不在于几何计算本身，而在于\textbf{在信息极度贫乏的条件下作出可检验的决策}。机器狗在任意时刻只知道已经执行过的动作及其返回结果：它不知道区域内有多少个干扰源，不知道各源的位置、接收半径，在第四问中还不知道源的类型与朝向。因此，任何"最优"都必须相对于"当前仍可能的世界"来定义，而不能相对于隐藏的真值。

这一观察决定了本文的统一技术路线。本文始终维护一个\textbf{可行集}——所有与已执行观测相容的世界的集合，并在该集合上取最坏值作为决策依据。具体地，位置层面的可行集是凸多面集 $P$（半平面之交），其半径 $\varrho(P)$ 直接度量"最坏情况下的定位误差"；当 $\varrho(P)\le20$ 时，无论真实源位于 $P$ 中何处，某个具体位置都能保证清除。这一"集合—最坏值"结构贯穿四问，只是集合的含义逐问扩展：问题一只有位置，问题二加入"接收半径不小于首次距离"的关系，问题三加入多源与未发现频道的存在性，问题四再把源类型与朝向（一个半平面形的可见性条件）纳入。

另一条贯穿全文的线索是\textbf{区分三类不同的约束}：由物理机制给出的硬约束（如 20 米清除距离、1000 米接收半径下界），由观测给出的信息约束（半平面），以及为求解效率而引入的工程参数（如横向分离阈值、评分权重）。第三类参数只在特定取值下经过数值检验，不能反过来被当作题目事实。本文在论述中明确标注每一处结论属于哪一类。

\subsection{问题一分析}

问题一要求给出计算定位区域直径的算法，并回答"以该直径为直径的圆能否覆盖该区域"。

其难点有二。第一，交会定位的经典做法是用斜率表示视线再求交，这在 $90^\circ$ 与 $270^\circ$ 附近会产生无穷斜率，在 $0^\circ$ 与 $360^\circ$ 的跨界处需要特殊分支。第二，"多边形定位区域"并不总是一个有界多边形：当观测方向接近平行、观测点相距很近或只做了很少几次观测时，可行域可能是无界区域，甚至是空集或一条线段。若不先判定区域状态就强行裁剪，会把无界性掩盖成"看起来有界"，从而使后续的直径与覆盖判断失去意义。

第三，也是最容易被忽略的一点：面积或直径小并不等于容易清除。清除要求存在一个位置到区域内\textbf{每一点}的距离都不超过 20 米，这是一个关于最小包围圆的条件，而不是关于直径的条件。因此问题一的正确输出不应止于直径，还须给出 $\varrho(P)$，这正是后两问真正需要的量。

基于以上分析，本文选择以向量半平面（叉积形式）表示角域，以线性规划判定区域状态，以顶点枚举计算直径，以最小包围圆处理覆盖问题，并用一个可真实构造的反例说明"以直径为直径的圆"一般不能覆盖。

\subsection{问题二分析}

问题二给出首次观测，要求选择第二个检测点并给出候选区域。它与问题一的联系在于：第二次观测的作用是缩小 $P$，评价标准自然落在更新后的区域上；与问题一的区别在于多出了一项硬约束——\textbf{第二个点必须还能收到信号}。如果第二点收不到信号，就得不到第二个示向度，交会定位无从进行。

这里的困难是：接收半径 $R$ 未知，真实距离 $r$ 也未知，二者取值范围不同（$R\in[1000,1500]$，而 $r$ 只要不超过 $R$ 即可）。若按最保守的方式要求"第二点到可行域内所有点的距离都不超过 1000 米"，会得到过于狭窄的候选域；若忽略接收条件只追求几何上好的交会角，则可能走到完全没有信号的死点上。

本文的关键观察是：首次返回\textbf{有效示向度}这一事件本身就携带信息——既然收到了信号，就必有 $r\le R$，而 $R$ 又有下界 1000。这一条看似简单的关系，使得可以构造出一个不依赖 $r$ 与 $R$ 具体取值的三圆盘区域，在其中任一点检测都保证能收到信号。在此区域内再排除与目标近似共线的退化位置，最后以最坏后验包围圆半径作为目标函数，即可把"较好的定位效果"这一模糊要求转化为一个可求解的 min-max 问题。

需要强调的是，题目并未指定精度与移动代价之间的权重。本文不编造单一权重强行宣布唯一最优点，而是给出随移动预算变化的策略族，并如实说明在较大预算下复杂搜索相对简单几何规则的增益很小。

\subsection{问题三分析}

问题三从单源、单步转入多源、整场，性质发生了变化。

首先是\textbf{目标函数变成整场总时间}。题目明确列出了计费项（移动、切换、检测、光学定位、清除），因此可以把整场写成式 \eqref{eq:time-ledger} 那样的可加账本。这一账本的意义在于：少走路、少检测、少切换、少失败尝试被放在同一时间尺度上比较，避免了仅优化其中一个分量的局部改进。由后续的官方演练可见，移动占 83\% 以上，因此路线是主要优化对象，但检测次数的收益也不能忽略。

其次出现\textbf{新的信息语义}：无信号。在全向源假设下，在某点收不到某频道的信号，说明该点与真实源的距离超过了接收半径，因此至少可以排除以该点为中心、半径 1000 米的圆盘。这一排除通常产生非凸集，本文用"保留圆盘外的凸外包"的方式处理，并进一步利用"同一源的接收半径固定"这一事实，把每一对正、负观测转化为一个半平面约束。

第三是\textbf{终止依据}。这是本题最容易出错、也最能体现建模水平的部分。题目说明源总数在 10 至 16 之间但具体个数未知，因此清除了 10 个或 14 个都不能结束；只有两种依据可以终止：其一，已清除 16 个，达到题目总数上限；其二，全部已发现源均已清除，且每个未发现频道都有全场覆盖证书。后者的实现方式是：把场地划为 128$\times$128 方格，只有当代方格的全部角点到某次实际无信号位置的距离都不超过 999.99 米时才排除该格（圆盘凸性保证整格被排除），当某频道的全部保留格均被排除时，才能证明该频道不存在。

第四是\textbf{发现与清除的联合调度}。已知目标的安全清除点与覆盖未发现频道的扫描点共同构成任务点集，需要在一个滚动时域内排序；每次实际观测后重新规划，只执行第一项。此外，当定位区域的半径已经足够小时，还可以考虑直接做一次光学试探以节约时间——这一决策同样要计入失败的可能与代价。

\subsection{问题四分析}

问题四引入定向源，使第三问的两项核心机制同时失效，必须逐项修正。

第一项失效在于\textbf{无信号的含义}。在全向情形下，某点无信号意味着"距离太远"；而在定向情形下，还可能是"源就在近处但正背向该点"，因为定向源只在朝向两侧各 $90^\circ$ 内辐射。因此第三问中"排除 1000 米圆盘"的操作在第四问中是错误的，会直接删掉真实源所在的位置。相应地，第三问把全向接收保证（三圆盘区域）用于第四问也是不允许的。

第二项失效在于\textbf{覆盖证书}。第三问的终止依据依赖"足以覆盖全场的扫描点集合"，其有效性建立在全向接收之上；对定向源，同样的站点布局无法保证发现任意朝向的源。一个极端的例子是：源位于 $(1800,0)$、朝向为 $(1,0)$，则严格位于场地内部的点全部处于其发射半平面的后方，仅在场内扫描没有任何发现保证。这也解释了为什么允许机器狗走出 1800 米圆域这一点是必要的——外侧扫描属于覆盖逻辑的一部分，而非数值越界。

本文的应对是给出两项\textbf{与朝向无关}的证书。其一是三角网证书：用边长 990 米的等边三角网覆盖场地，若某频道在某个三角形的三个顶点都实际测得无信号，那么由凸组合可知，任意朝向下至少有一个顶点位于源的闭发射半平面内，且该顶点到源的距离不超过 990 米、小于接收半径下界 1000 米，故该顶点必能收到信号，矛盾。其二是更一般的凸包排除条件：对任一凸子区域 $Q$，只要存在一组到 $Q$ 内所有点距离都不超过 1000 米的实际无信号站点，且 $Q$ 包含于这些站点的凸包，则 $Q$ 内不存在任意朝向的定向源。后者允许 $Q$ 的不同顶点由不同的站点组合证明，从而把原先固定的 31 点网格替换为更少的站点，这正是第四问效率提升的主要来源。

至于正观测条件下的定位、包围圆与清除判据，本文原样继承第三问——因为光学精确定位与清除是否成功只与距离有关，与源的朝向无关。这一"可继承/不可继承"的清单是第四问建模的起点，也是本文在论述中刻意分开写清的部分。
